\section*{Примечания о чтениях}
\subsection*{\texttt{GELFAND-P001465}}
\par\noindent\leavevmode\phantomsection\label{GELFAND-0159-U001}\label{GELFAND-P001465:eq18-composition}\textbf{\texttt{GELFAND-0159-U001}} (eq18-composition).\nopagebreak[4]\par
В (18) видны блок BA, затем соединённые блоки B и A, затем сумма и два раздельных блока. Штрихи у верхних j в среднем рисунке близки к двойным; правый нижний индекс у A пересекается с нитью и похож на M, хотя ниже стоит m\_L.\par
Кандидаты: Сверху B: j\_1'', j\_2'', j\_M; снизу A: j\_1, j\_2, j\_M; У отдельных верхних j один штрих; последний нижний индекс у A: L. Чтение не установлено в дипломатическом исходнике.\par
Оригинал 2408×3480 px. Вся область (65,100,2040,975); контекстные увеличения 3×: (65,100,715,735), (865,155,1365,975), (1370,220,2040,970); детали 8×: (180,180,600,380), (930,185,1400,340), (1165,730,1320,865). Индексы не согласованы по ожидаемому правилу композиции; описание сохраняет видимые обозначения, включая несовпадение M/L. Чтение верхних окончаний у B затруднено наложением строки заголовка.\par
Повторная сверка: Повторно просмотрены исходник и увеличения 3×, 5× и 10×. Исходные наблюдения, варианты чтения, зачёркивания и неуверенности сохранены. Математические проверки отделены от буквального чтения; частичные рисунки не утверждают неустановленных соединений или направлений.\par\medskip
\par\noindent\leavevmode\phantomsection\label{GELFAND-0159-U002}\label{GELFAND-P001465:eq20-crossing}\textbf{\texttt{GELFAND-0159-U002}} (eq20-crossing).\nopagebreak[4]\par
У нижнего правого конца непрерывной диагонали в рисунке (20) короткий боковой загнутый штрих примыкает к заострённому концу. Второе крыло стрелки отдельно не различается. Подписи и разрыв второй нити читаются.\par
Кандидаты: стрелка вниз вправо; обведённый заострённый конец без стрелки. Чтение не установлено в дипломатическом исходнике.\par
Оригинал 2408×3480 px; область (220,1250,515,1495), контекст (60,950,1810,1515). Осмотрено в оригинале и увеличении 3×, дополнительно независимо сверено по увеличению. Направленность нити не восстановлена по последующим рисункам. В (20') такого бокового штриха у окончания нет.\par
Повторная сверка: Повторно просмотрены исходник и увеличения 3×, 5× и 10×. Исходные наблюдения, варианты чтения, зачёркивания и неуверенности сохранены. Математические проверки отделены от буквального чтения; частичные рисунки не утверждают неустановленных соединений или направлений.\par\medskip
\par\noindent\leavevmode\phantomsection\label{GELFAND-0159-U003}\label{GELFAND-P001465:eq20-v-indices}\textbf{\texttt{GELFAND-0159-U003}} (eq20-v-indices).\nopagebreak[4]\par
В R\textasciicircum{}\{j\_1 j\_2\} второй индекс выглядит как 2. Во всех четырёх V в той же строке индекс при j читается как 1; отдельный просмотр подтвердил одинаковый прямой штрих с нижней ножкой.\par
Кандидаты: R\textasciicircum{}\{j\_1j\_2\}:V\textasciicircum{}\{j\_1\}\textbackslash{}otimes V\textasciicircum{}\{j\_1\}\textbackslash{}longrightarrow V\textasciicircum{}\{j\_1\}\textbackslash{}otimes V\textasciicircum{}\{j\_1\}. Чтение не установлено в дипломатическом исходнике.\par
Оригинал 2408×3480 px; область (700,1220,1790,1400), увеличение 8×; общий контекст (60,950,1810,1515), увеличение 3×. Повторение j\_1 сохранено буквально. Соседние рисунки имеют цвета j\_1,j\_2, но это не основание подменять индексы в формуле. Нотационное согласование остаётся на проверку; исправленный вариант не утверждается.\par
Повторная сверка: Повторно просмотрены исходник и увеличения 3×, 5× и 10×. Исходные наблюдения, варианты чтения, зачёркивания и неуверенности сохранены. Математические проверки отделены от буквального чтения; частичные рисунки не утверждают неустановленных соединений или направлений.\par\medskip
\subsection*{\texttt{GELFAND-P001466}}
\par\noindent\leavevmode\phantomsection\label{GELFAND-0159-U004}\label{GELFAND-P001466:eq25-cancelled-exponent}\textbf{\texttt{GELFAND-0159-U004}} (eq25-cancelled-exponent).\nopagebreak[4]\par
В (25) показатель у (-1) начинается с m, справа — густой перечёркнутый знак с несколькими слившимися штрихами. Нельзя надёжно отделить исходный индекс от перечёркивания. Показатель q справа читается −m\_1/2.\par
Кандидаты: m\_1; m с неразборчивым перечёркнутым знаком справа. Чтение не установлено в дипломатическом исходнике.\par
Оригинал 2408×3480 px; область (1015,1265,1365,1375), увеличение 10×; контекст (105,795,1670,1540), 3×. Кандидат m\_1 в теле помечен как неуверенный; зачёркнутый знак отдельно не восстановлен. Никакой дополнительный множитель или цифра не добавлены. В (24) показатель q видимо m\_1, в (25) −m\_1/2; оба сохранены без унификации.\par
Повторная сверка: Повторно просмотрены исходник и увеличения 3×, 5× и 10×. Исходные наблюдения, варианты чтения, зачёркивания и неуверенности сохранены. Математические проверки отделены от буквального чтения; частичные рисунки не утверждают неустановленных соединений или направлений.\par\medskip
\par\noindent\leavevmode\phantomsection\label{GELFAND-0159-U005}\label{GELFAND-P001466:eq26-vertex-arrow}\textbf{\texttt{GELFAND-0159-U005}} (eq26-vertex-arrow).\nopagebreak[4]\par
В правом рисунке (26) на ветви j\_2 ниже и правее вершины стрелка покрыта густыми повторными штрихами. На верхней правой и нижней левой стороне выступают короткие части знака, допускающие направление к вершине.\par
Кандидаты: стрелка j\_2 к вершине, вверх влево; направление стрелки j\_2 неразборчиво. Чтение не установлено в дипломатическом исходнике.\par
Оригинал 2408×3480 px; вся область (140,1660,1365,2030); контекст (140,1550,1555,2465), увеличение 3×; деталь (1190,1790,1320,1925), 10×. Читаемые ветви и стрелки обеих частей равенства перечислены в SourceFigure. Направление неоднозначной стрелки не задано по алгебраическому ожиданию.\par
Повторная сверка: Повторно просмотрены исходник и увеличения 3×, 5× и 10×. Исходные наблюдения, варианты чтения, зачёркивания и неуверенности сохранены. Математические проверки отделены от буквального чтения; частичные рисунки не утверждают неустановленных соединений или направлений.\par\medskip
\par\noindent\leavevmode\phantomsection\label{GELFAND-0159-U006}\label{GELFAND-P001466:eq28-coefficient-j}\textbf{\texttt{GELFAND-0159-U006}} (eq28-coefficient-j).\nopagebreak[4]\par
В (28) обе выходящие нижние нити обозначены j\_3. В коэффициенте показатель (-1) заканчивается на −j, а числитель показателя q начинается с j(j+1); нижний индекс 3 при этих j не виден. Знаменатель показателя q — 4.\par
Кандидаты: (-1)\textasciicircum{}\{j\_1+j\_2-j\}q\textasciicircum{}\{(j(j+1)-j\_1(j\_1+1)-j\_2(j\_2+1))/4\}. Чтение не установлено в дипломатическом исходнике.\par
Оригинал 2408×3480 px; область (140,2495,2040,3020) внутри просмотренного увеличения (140,2495,2040,3300), 3×. Буквальное j в коэффициенте и j\_3 на рисунках сохранены; возможное отождествление этих обозначений не внесено в источник как установленное равенство.\par
Повторная сверка: Повторно просмотрены исходник и увеличения 3×, 5× и 10×. Исходные наблюдения, варианты чтения, зачёркивания и неуверенности сохранены. Математические проверки отделены от буквального чтения; частичные рисунки не утверждают неустановленных соединений или направлений.\par\medskip
\subsection*{\texttt{GELFAND-P001467}}
\par\noindent\leavevmode\phantomsection\label{GELFAND-0159-U201}\label{GELFAND-P001467:x58-y247-x646-y638}\textbf{\texttt{GELFAND-0159-U201}} (x58-y247-x646-y638).\nopagebreak[4]\par
Рисунок (29): слева стрелка вверх с подписью t, справа криволинейная петля; верхний и нижний изгибы и два места в верхней средней части обведены овалами. Над левым средним местом стоит m\_i, справа — j. Видны пересечения, стрелочные штрихи и разрывы линий справа.\par
Кандидаты: не записаны. Чтение не установлено в дипломатическом исходнике.\par
Координаты исходного изображения 2408×3480: x=58–646, y=247–638, начало в левом верхнем углу. Проверено на полном изображении и при увеличении 3×. Из-за наложения штрихов точный ход ветвей в пересечениях и все направления стрелок не установлены; дано описание видимого рисунка без выдуманной топологии.\par
Повторная сверка: Повторно просмотрены исходник и увеличения 3×, 5× и 10×. Исходные наблюдения, варианты чтения, зачёркивания и неуверенности сохранены. Математические проверки отделены от буквального чтения; частичные рисунки не утверждают неустановленных соединений или направлений.\par\medskip
\par\noindent\leavevmode\phantomsection\label{GELFAND-0159-U202}\label{GELFAND-P001467:x656-y1355-x935-y1489}\textbf{\texttt{GELFAND-0159-U202}} (x656-y1355-x935-y1489).\nopagebreak[4]\par
В (30) между знаком произведения и скобками (\{m\}) стоит небольшой v-образный знак; от него идёт указатель к подписи о вычисленных выше матричных элементах.\par
Кандидаты: v; \textbackslash{}nu. Чтение не установлено в дипломатическом исходнике.\par
Координаты исходного изображения 2408×3480: x=656–935, y=1355–1489. Проверено при увеличениях 3× и 10×; форма не позволяет уверенно различить латинскую v и греческую nu. В формуле дан кандидат v с явным маркером неопределённости. На соседнем предшествующем листе определение этого знака не найдено.\par
Повторная сверка: Повторно просмотрены исходник и увеличения 3×, 5× и 10×. Исходные наблюдения, варианты чтения, зачёркивания и неуверенности сохранены. Математические проверки отделены от буквального чтения; частичные рисунки не утверждают неустановленных соединений или направлений.\par\medskip
\par\noindent\leavevmode\phantomsection\label{GELFAND-0159-U203}\label{GELFAND-P001467:x1935-y1670-x2000-y1726}\textbf{\texttt{GELFAND-0159-U203}} (x1935-y1670-x2000-y1726).\nopagebreak[4]\par
В (31) непосредственно после N\_-\textasciicircum{}i(D\_L), в верхнем правом конце показателя степени q, имеется небольшой зачёркнутый надстрочный знак.\par
Кандидаты: не записаны. Чтение не установлено в дипломатическом исходнике.\par
Координаты исходного изображения 2408×3480: x=1935–2000, y=1670–1726. Просмотрено в составе формулы при увеличениях 3× и 5× и отдельно при увеличении 10×; исходный знак под плотными пересекающимися штрихами не читается. Зачёркивание сохранено описательно.\par
Повторная сверка: Повторно просмотрены исходник и увеличения 3×, 5× и 10×. Исходные наблюдения, варианты чтения, зачёркивания и неуверенности сохранены. Математические проверки отделены от буквального чтения; частичные рисунки не утверждают неустановленных соединений или направлений.\par\medskip
\par\noindent\leavevmode\phantomsection\label{GELFAND-0159-U204}\label{GELFAND-P001467:x520-y2610-x734-y2730}\textbf{\texttt{GELFAND-0159-U204}} (x520-y2610-x734-y2730).\nopagebreak[4]\par
В строке «Можно … строить инв. узлов» между «Можно» и «строить» зачёркнуто короткое слово.\par
Кандидаты: не записаны. Чтение не установлено в дипломатическом исходнике.\par
Координаты исходного изображения 2408×3480: x=520–734, y=2610–2730. После просмотров при увеличениях 3× и 10× чтение скрытого слова не установлено. Остальная строка транскрибирована.\par
Повторная сверка: Повторно просмотрены исходник и увеличения 3×, 5× и 10×. Исходные наблюдения, варианты чтения, зачёркивания и неуверенности сохранены. Математические проверки отделены от буквального чтения; частичные рисунки не утверждают неустановленных соединений или направлений.\par\medskip
\subsection*{\texttt{GELFAND-P001468}}
\par\noindent\leavevmode\phantomsection\label{GELFAND-0159-U205}\label{GELFAND-P001468:x1260-y33-x1400-y148}\textbf{\texttt{GELFAND-0159-U205}} (x1260-y33-x1400-y148).\nopagebreak[4]\par
В верхней строке между подчёркнутой фамилией «Ушверидзе» и именем «Александр» находится густо зачёркнутая короткая запись.\par
Кандидаты: не записаны. Чтение не установлено в дипломатическом исходнике.\par
Координаты исходного изображения 2408×3480: x=1260–1400, y=33–148. Просмотрено при увеличениях 3× и 10×; буквы или инициалы под штрихами не прочитаны. Незачёркнутые имя и отчество сохранены.\par
Повторная сверка: Повторно просмотрены исходник и увеличения 3×, 5× и 10×. Исходные наблюдения, варианты чтения, зачёркивания и неуверенности сохранены. Математические проверки отделены от буквального чтения; частичные рисунки не утверждают неустановленных соединений или направлений.\par\medskip
\par\noindent\leavevmode\phantomsection\label{GELFAND-0159-U206}\label{GELFAND-P001468:x2040-y40-x2310-y150}\textbf{\texttt{GELFAND-0159-U206}} (x2040-y40-x2310-y150).\nopagebreak[4]\par
Справа в верхней строке, после «Григорьевич.», имеется запись, перечёркнутая несколькими почти горизонтальными штрихами.\par
Кандидаты: не записаны. Чтение не установлено в дипломатическом исходнике.\par
Координаты исходного изображения 2408×3480: x=2040–2310, y=40–150. Проверено на полном изображении и при увеличениях 3× и 10×. Исходные знаки сливаются с отменяющими штрихами; чтение не установлено.\par
Повторная сверка: Повторно просмотрены исходник и увеличения 3×, 5× и 10×. Исходные наблюдения, варианты чтения, зачёркивания и неуверенности сохранены. Математические проверки отделены от буквального чтения; частичные рисунки не утверждают неустановленных соединений или направлений.\par\medskip
\par\noindent\leavevmode\phantomsection\label{GELFAND-0159-U207}\label{GELFAND-P001468:x1000-y1240-x1248-y1328}\textbf{\texttt{GELFAND-0159-U207}} (x1000-y1240-x1248-y1328).\nopagebreak[4]\par
Над строкой после (4) вставлено «условиях на Δ …»; последнее сокращённое слово вставки начинается на «нахо» и заканчивается перед точкой.\par
Кандидаты: нахожд.. Чтение не установлено в дипломатическом исходнике.\par
Координаты исходного изображения 2408×3480: x=1000–1248, y=1240–1328. Проверено при увеличениях 5× и 10×. Последние соединённые буквы не вполне различимы; «нахожд.» дано как кандидат, без развёртывания сокращения. Вставка и её место отмечены в тексте.\par
Повторная сверка: Повторно просмотрены исходник и увеличения 3×, 5× и 10×. Исходные наблюдения, варианты чтения, зачёркивания и неуверенности сохранены. Математические проверки отделены от буквального чтения; частичные рисунки не утверждают неустановленных соединений или направлений.\par\medskip
\par\noindent\leavevmode\phantomsection\label{GELFAND-0159-U208}\label{GELFAND-P001468:x435-y2450-x570-y2578}\textbf{\texttt{GELFAND-0159-U208}} (x435-y2450-x570-y2578).\nopagebreak[4]\par
В пункте (6), 3), перед незачёркнутым равенством Ω(λ)−Ω(−λ)=0 имеется зачёркнутая короткая формула; в её левой части различим треугольный знак, похожий на Δ.\par
Кандидаты: не записаны. Чтение не установлено в дипломатическом исходнике.\par
Координаты исходного изображения 2408×3480: x=435–570, y=2450–2578. Проверено при увеличении 10×. Скобкообразные и буквенные штрихи перекрыты зачёркиванием, поэтому полный прежний фрагмент не восстанавливается.\par
Повторная сверка: Повторно просмотрены исходник и увеличения 3×, 5× и 10×. Исходные наблюдения, варианты чтения, зачёркивания и неуверенности сохранены. Математические проверки отделены от буквального чтения; частичные рисунки не утверждают неустановленных соединений или направлений.\par\medskip
\par\noindent\leavevmode\phantomsection\label{GELFAND-0159-U209}\label{GELFAND-P001468:x1335-y2810-x1494-y2960}\textbf{\texttt{GELFAND-0159-U209}} (x1335-y2810-x1494-y2960).\nopagebreak[4]\par
Справа от первой пары формул в (7), после Ω(λ)=0;, зачёркнут небольшой фрагмент; различим треугольный знак, похожий на Δ.\par
Кандидаты: не записаны. Чтение не установлено в дипломатическом исходнике.\par
Координаты исходного изображения 2408×3480: x=1335–1494, y=2810–2960. Просмотрено при увеличениях 3× и 10×. Остальные знаки под плотными штрихами не установлены; математически ожидаемая запись не подставлялась.\par
Повторная сверка: Повторно просмотрены исходник и увеличения 3×, 5× и 10×. Исходные наблюдения, варианты чтения, зачёркивания и неуверенности сохранены. Математические проверки отделены от буквального чтения; частичные рисунки не утверждают неустановленных соединений или направлений.\par\medskip
\par\noindent\leavevmode\phantomsection\label{GELFAND-0159-U213}\label{GELFAND-P001468:x350-y1390-x1615-y1675}\textbf{\texttt{GELFAND-0159-U213}} (x350-y1390-x1615-y1675).\nopagebreak[4]\par
Первая строка (5) написана как sum\_\{i=1\}\textasciicircum{}M prime Delta(lambda\_i-lambda\_k) + sum\_\{alpha=1\}\textasciicircum{}N b\_alpha(lambda\_i-a\_alpha)=0; под первым суммированием указано «сумма по i≠k». После b\_alpha непосредственно открывается скобка, без знака Delta; индекс у lambda во второй скобке именно i.\par
Кандидаты: не записаны. Чтение не установлено в дипломатическом исходнике.\par
Координаты исходного изображения 2408×3480: x=350–1615, y=1390–1675. Проверено на полном изображении и по 3× увеличению формулы (5). Отсутствие Delta и сочетание суммируемого i с индексом i во втором члене требуют математической сверки с (1)–(4). Это замечание о возможной ошибке источника; все видимые символы, индексы и примечание сохранены, исправленное уравнение не подставлялось.\par
Повторная сверка: Повторно просмотрены исходник и увеличения 3×, 5× и 10×. Исходные наблюдения, варианты чтения, зачёркивания и неуверенности сохранены. Математические проверки отделены от буквального чтения; частичные рисунки не утверждают неустановленных соединений или направлений.\par\medskip
\subsection*{\texttt{GELFAND-P001469}}
\par\noindent\leavevmode\phantomsection\label{GELFAND-0159-U210}\label{GELFAND-P001469:x825-y1238-x1030-y1340}\textbf{\texttt{GELFAND-0159-U210}} (x825-y1238-x1030-y1340).\nopagebreak[4]\par
Во вставке «в случае … (7)\_1» над пунктом В перед незачёркнутой ссылкой (7)\_1 видна зачёркнутая запись.\par
Кандидаты: не записаны. Чтение не установлено в дипломатическом исходнике.\par
Координаты исходного изображения 2408×3480: x=825–1030, y=1238–1340. Просмотрено при увеличениях 3× и 10×; прежний номер или формульный знак под перекрёстными штрихами не различим. Вставка и видимая ссылка сохранены.\par
Повторная сверка: Повторно просмотрены исходник и увеличения 3×, 5× и 10×. Исходные наблюдения, варианты чтения, зачёркивания и неуверенности сохранены. Математические проверки отделены от буквального чтения; частичные рисунки не утверждают неустановленных соединений или направлений.\par\medskip
\par\noindent\leavevmode\phantomsection\label{GELFAND-0159-U211}\label{GELFAND-P001469:x800-y1425-x1005-y1558}\textbf{\texttt{GELFAND-0159-U211}} (x800-y1425-x1005-y1558).\nopagebreak[4]\par
В пункте В после «подвижные с» и перед λ\_i зачёркнуто слово; над ним написано «коорд.».\par
Кандидаты: не записаны. Чтение не установлено в дипломатическом исходнике.\par
Координаты исходного изображения 2408×3480: x=800–1005, y=1425–1558. Проверено при увеличениях 3× и 10×. Исправленная запись «коорд.» читается, но исходное зачёркнутое слово не установлено; оба слоя отмечены в теле.\par
Повторная сверка: Повторно просмотрены исходник и увеличения 3×, 5× и 10×. Исходные наблюдения, варианты чтения, зачёркивания и неуверенности сохранены. Математические проверки отделены от буквального чтения; частичные рисунки не утверждают неустановленных соединений или направлений.\par\medskip
\par\noindent\leavevmode\phantomsection\label{GELFAND-0159-U212}\label{GELFAND-P001469:x500-y2640-x570-y2745}\textbf{\texttt{GELFAND-0159-U212}} (x500-y2640-x570-y2745).\nopagebreak[4]\par
В (10), под строкой «В кулоновской сист.», непосредственно перед λ\_i зачёркнут отдельный знак.\par
Кандидаты: не записаны. Чтение не установлено в дипломатическом исходнике.\par
Координаты исходного изображения 2408×3480: x=500–570, y=2640–2745. Просмотрено при увеличении 10×. Наклонные штрихи отмены перекрывают исходный знак; чтение не установлено.\par
Повторная сверка: Повторно просмотрены исходник и увеличения 3×, 5× и 10×. Исходные наблюдения, варианты чтения, зачёркивания и неуверенности сохранены. Математические проверки отделены от буквального чтения; частичные рисунки не утверждают неустановленных соединений или направлений.\par\medskip
\par\noindent\leavevmode\phantomsection\label{GELFAND-0159-U214}\label{GELFAND-P001469:x438-y860-x2110-y1115}\textbf{\texttt{GELFAND-0159-U214}} (x438-y860-x2110-y1115).\nopagebreak[4]\par
В (9) написано Delta\_0(b lambda)+a lambda — «общее решение сист. (6)»; внешнего множителя перед Delta\_0 нет. Следующая строка: «Delta\_0 — одна из ф-й примера (6)».\par
Кандидаты: не записаны. Чтение не установлено в дипломатическом исходнике.\par
Координаты исходного изображения 2408×3480: x=438–2110, y=860–1115. Проверено на полном изображении; начальная формула повторно просмотрена при увеличении 5×. Вид формулы общего решения требует математической сверки с (6). Ссылка «примера (6)» также оставлена буквально: примеры на предшествующем кадре P001468 имеют номер (7), тогда как (6) обозначает свойства Delta и Omega. Никаких множителей и исправленных номеров не добавлено.\par
Повторная сверка: Повторно просмотрены исходник и увеличения 3×, 5× и 10×. Исходные наблюдения, варианты чтения, зачёркивания и неуверенности сохранены. Математические проверки отделены от буквального чтения; частичные рисунки не утверждают неустановленных соединений или направлений.\par\medskip
\subsection*{\texttt{GELFAND-P001470}}
\par\noindent\leavevmode\phantomsection\label{GELFAND-0159-U401}\label{GELFAND-P001470:top-left-label}\textbf{\texttt{GELFAND-0159-U401}} (top-left-label).\nopagebreak[4]\par
На верхнем левом поле видна густо зачёркнутая помета; отдельный первоначальный номер или знак уверенно не читается.\par
Кандидаты: не записаны. Чтение не установлено в дипломатическом исходнике.\par
Исходный PNG 2408×3480; область x=130–240, y=80–191 пикселей. Просмотрен весь оригинал, затем область увеличена до 8×. В SourcePage оставлено unlabelled; номер соседней страницы не перенесён на этот кадр.\par
Повторная сверка: Повторно просмотрены исходник и увеличения 3×, 5× и 10×. Исходные наблюдения, варианты чтения, зачёркивания и неуверенности сохранены. Математические проверки отделены от буквального чтения; частичные рисунки не утверждают неустановленных соединений или направлений.\par\medskip
\par\noindent\leavevmode\phantomsection\label{GELFAND-0159-U402}\label{GELFAND-P001470:top-line-cancellation}\textbf{\texttt{GELFAND-0159-U402}} (top-line-cancellation).\nopagebreak[4]\par
После «тогда» в мелкой верхней строке имеется одно короткое слово, перекрытое несколькими горизонтальными штрихами.\par
Кандидаты: не записаны. Чтение не установлено в дипломатическом исходнике.\par
Исходный PNG 2408×3480; область x=1110–1300, y=30–112 пикселей. После просмотра увеличения 10× остатки букв не дают уверенного чтения зачёркнутого слова.\par
Повторная сверка: Повторно просмотрены исходник и увеличения 3×, 5× и 10×. Исходные наблюдения, варианты чтения, зачёркивания и неуверенности сохранены. Математические проверки отделены от буквального чтения; частичные рисунки не утверждают неустановленных соединений или направлений.\par\medskip
\par\noindent\leavevmode\phantomsection\label{GELFAND-0159-U403}\label{GELFAND-P001470:top-right-cancellation}\textbf{\texttt{GELFAND-0159-U403}} (top-right-cancellation).\nopagebreak[4]\par
Между «соб. знач.» и окончательным H имеется зачёркнутая группа крупных наклонных штрихов.\par
Кандидаты: не записаны. Чтение не установлено в дипломатическом исходнике.\par
Исходный PNG 2408×3480; область x=1950–2130, y=70–195 пикселей. Просмотрено при 10×; зачёркивание скрывает исходные знаки. Окончательное H справа читается и передано отдельно.\par
Повторная сверка: Повторно просмотрены исходник и увеличения 3×, 5× и 10×. Исходные наблюдения, варианты чтения, зачёркивания и неуверенности сохранены. Математические проверки отделены от буквального чтения; частичные рисунки не утверждают неустановленных соединений или направлений.\par\medskip
\par\noindent\leavevmode\phantomsection\label{GELFAND-0159-U404}\label{GELFAND-P001470:top-eigenvectors-verb}\textbf{\texttt{GELFAND-0159-U404}} (top-eigenvectors-verb).\nopagebreak[4]\par
В скобках после «и» перед «соб. вект.» рукописное слово имеет начальную петлю с нижним выносом и слитные средние буквы.\par
Кандидаты: знаем. Чтение не установлено в дипломатическом исходнике.\par
Исходный PNG 2408×3480; область x=388–670, y=166–263 пикселей. Повторно просмотрены соседние строки и увеличения 6× и 10×. Чтение «знаем» оставлено кандидатом, без восстановления по смыслу.\par
Повторная сверка: Повторно просмотрены исходник и увеличения 3×, 5× и 10×. Исходные наблюдения, варианты чтения, зачёркивания и неуверенности сохранены. Математические проверки отделены от буквального чтения; частичные рисунки не утверждают неустановленных соединений или направлений.\par\medskip
\par\noindent\leavevmode\phantomsection\label{GELFAND-0159-U405}\label{GELFAND-P001470:lower-number-cancellation}\textbf{\texttt{GELFAND-0159-U405}} (lower-number-cancellation).\nopagebreak[4]\par
После «числах» перед c\_α зачёркнута небольшая группа штрихов.\par
Кандидаты: не записаны. Чтение не установлено в дипломатическом исходнике.\par
Исходный PNG 2408×3480; область x=557–666, y=2350–2476 пикселей. Просмотрено при 10×. Числа c\_α, λ\_i справа от зачёркивания переданы отдельно; зачёркнутый знак не восстанавливался.\par
Повторная сверка: Повторно просмотрены исходник и увеличения 3×, 5× и 10×. Исходные наблюдения, варианты чтения, зачёркивания и неуверенности сохранены. Математические проверки отделены от буквального чтения; частичные рисунки не утверждают неустановленных соединений или направлений.\par\medskip
\par\noindent\leavevmode\phantomsection\label{GELFAND-0159-U406}\label{GELFAND-P001470:lower-equation-cancellation}\textbf{\texttt{GELFAND-0159-U406}} (lower-equation-cancellation).\nopagebreak[4]\par
После «Тогда (13)» перед знаком подобия имеется короткая зачёркнутая запись с нижним выносным штрихом.\par
Кандидаты: ур.. Чтение не установлено в дипломатическом исходнике.\par
Исходный PNG 2408×3480; область x=790–952, y=2650–2780 пикселей. Просмотрено при 10×. Сохранившийся вынос допускает сокращение «ур.», но густое перечёркивание не позволяет утверждать чтение.\par
Повторная сверка: Повторно просмотрены исходник и увеличения 3×, 5× и 10×. Исходные наблюдения, варианты чтения, зачёркивания и неуверенности сохранены. Математические проверки отделены от буквального чтения; частичные рисунки не утверждают неустановленных соединений или направлений.\par\medskip
\subsection*{\texttt{GELFAND-P001471}}
\par\noindent\leavevmode\phantomsection\label{GELFAND-0159-U407}\label{GELFAND-P001471:top-initial-cancellation}\textbf{\texttt{GELFAND-0159-U407}} (top-initial-cancellation).\nopagebreak[4]\par
В строке (17) перед «В случае» находится густо зачёркнутый отдельный знак.\par
Кандидаты: не записаны. Чтение не установлено в дипломатическом исходнике.\par
Исходный PNG 2408×3480; область x=310–412, y=66–182 пикселей. Просмотрено при 10×; первоначальная буква или символ не различается уверенно.\par
Повторная сверка: Повторно просмотрены исходник и увеличения 3×, 5× и 10×. Исходные наблюдения, варианты чтения, зачёркивания и неуверенности сохранены. Математические проверки отделены от буквального чтения; частичные рисунки не утверждают неустановленных соединений или направлений.\par\medskip
\par\noindent\leavevmode\phantomsection\label{GELFAND-0159-U408}\label{GELFAND-P001471:top-line-cancellation}\textbf{\texttt{GELFAND-0159-U408}} (top-line-cancellation).\nopagebreak[4]\par
Под словами «В случае (7)\_1» отдельная короткая строка зачёркнута несколькими горизонтальными линиями; видны части «лин.», «ф-я» и последней λ.\par
Кандидаты: лин. ф-я от λ. Чтение не установлено в дипломатическом исходнике.\par
Исходный PNG 2408×3480; область x=445–925, y=195–285 пикселей. Строка осмотрена при 6× и 10×. Нижнее незачёркнутое пояснение также содержит «Ξ — лин. ф-я от λ», однако повтор не использован как доказательство скрытых букв в зачёркнутой строке.\par
Повторная сверка: Повторно просмотрены исходник и увеличения 3×, 5× и 10×. Исходные наблюдения, варианты чтения, зачёркивания и неуверенности сохранены. Математические проверки отделены от буквального чтения; частичные рисунки не утверждают неустановленных соединений или направлений.\par\medskip
\par\noindent\leavevmode\phantomsection\label{GELFAND-0159-U409}\label{GELFAND-P001471:middle-potential-cancellation}\textbf{\texttt{GELFAND-0159-U409}} (middle-potential-cancellation).\nopagebreak[4]\par
В строке (18) перед «Потенц. V» зачёркнуто слово; часть первого знака похожа на Е, далее видны несколько петель.\par
Кандидаты: Если. Чтение не установлено в дипломатическом исходнике.\par
Исходный PNG 2408×3480; область x=372–630, y=1565–1715 пикселей. Просмотрено при 8×. Чтение первоначального слова остаётся неопределённым, поскольку буквы пересечены широкими наклонными штрихами.\par
Повторная сверка: Повторно просмотрены исходник и увеличения 3×, 5× и 10×. Исходные наблюдения, варианты чтения, зачёркивания и неуверенности сохранены. Математические проверки отделены от буквального чтения; частичные рисунки не утверждают неустановленных соединений или направлений.\par\medskip
\par\noindent\leavevmode\phantomsection\label{GELFAND-0159-U413}\label{GELFAND-P001471:diagram20-faint-strokes}\textbf{\texttt{GELFAND-0159-U413}} (diagram20-faint-strokes).\nopagebreak[4]\par
У левого конца горизонтали в схеме (20) имеются утолщение и наклонные штрихи, допускающие чтение как наконечника стрелки; ниже этого конца и тёмной точки со знаком минус видны слабые прерывистые штрихи вниз влево. Остальные жирные элементы схемы читаются отчётливо.\par
Кандидаты: стрелка влево на левом конце горизонтали. Чтение не установлено в дипломатическом исходнике.\par
Исходный PNG 2408×3480; область x=340–540, y=2630–2760 пикселей. Осмотрено при 6× и 8×. Читаемые элементы воспроизведены в TikZ; неоднозначный левый конец и слабые штрихи описаны SourceFigure без уверенного назначения направления или физических связей. Отдельная читаемая стрелка на правой части линии направлена влево.\par
Повторная сверка: Повторно просмотрены исходник и увеличения 3×, 5× и 10×. Исходные наблюдения, варианты чтения, зачёркивания и неуверенности сохранены. Математические проверки отделены от буквального чтения; частичные рисунки не утверждают неустановленных соединений или направлений.\par\medskip
\par\noindent\leavevmode\phantomsection\label{GELFAND-0159-U414}\label{GELFAND-P001471:margin-below20}\textbf{\texttt{GELFAND-0159-U414}} (margin-below20).\nopagebreak[4]\par
На левом поле под номером (20) имеется короткий отдельный слегка изогнутый наклонный штрих.\par
Кандидаты: не записаны. Чтение не установлено в дипломатическом исходнике.\par
Исходный PNG 2408×3480; область x=70–108, y=2750–2799 пикселей. Осмотрено при 8×. Отдельный штрих учтён без догадки о буквенном или пунктуационном значении.\par
Повторная сверка: Повторно просмотрены исходник и увеличения 3×, 5× и 10×. Исходные наблюдения, варианты чтения, зачёркивания и неуверенности сохранены. Математические проверки отделены от буквального чтения; частичные рисунки не утверждают неустановленных соединений или направлений.\par\medskip
\subsection*{\texttt{GELFAND-P001472}}
\par\noindent\leavevmode\phantomsection\label{GELFAND-0159-U410}\label{GELFAND-P001472:upper-atiya-name}\textbf{\texttt{GELFAND-0159-U410}} (upper-atiya-name).\nopagebreak[4]\par
В подзаголовке «о раб.» короткая фамилия начинается «Ати»; конечные слитные штрихи не позволяют уверенно решить, есть ли ещё одна и.\par
Кандидаты: Ати; Атии. Чтение не установлено в дипломатическом исходнике.\par
Исходный PNG 2408×3480; область x=627–878, y=470–560 пикселей. Повторно просмотрено при 8×. Каталожные сведения не использованы для расширения или исправления видимой формы фамилии.\par
Повторная сверка: Повторно просмотрены исходник и увеличения 3×, 5× и 10×. Исходные наблюдения, варианты чтения, зачёркивания и неуверенности сохранены. Математические проверки отделены от буквального чтения; частичные рисунки не утверждают неустановленных соединений или направлений.\par\medskip
\par\noindent\leavevmode\phantomsection\label{GELFAND-0159-U411}\label{GELFAND-P001472:middle-flow-cancellation}\textbf{\texttt{GELFAND-0159-U411}} (middle-flow-cancellation).\nopagebreak[4]\par
После «Рассм.» на основной строке зачёркнуто слово; над ним вписано читаемое «траектории».\par
Кандидаты: не записаны. Чтение не установлено в дипломатическом исходнике.\par
Исходный PNG 2408×3480; область x=530–830, y=1872–1985 пикселей. Просмотрено при 8×. Верхняя вставка передана полностью; нижняя первоначальная запись скрыта несколькими штрихами и не восстановлена.\par
Повторная сверка: Повторно просмотрены исходник и увеличения 3×, 5× и 10×. Исходные наблюдения, варианты чтения, зачёркивания и неуверенности сохранены. Математические проверки отделены от буквального чтения; частичные рисунки не утверждают неустановленных соединений или направлений.\par\medskip
\par\noindent\leavevmode\phantomsection\label{GELFAND-0159-U412}\label{GELFAND-P001472:bottom-orientation-cancellation}\textbf{\texttt{GELFAND-0159-U412}} (bottom-orientation-cancellation).\nopagebreak[4]\par
После «ориентац.» зачёркнута запись с различимыми нижними знаками, похожими на индексы 1 и 2; над ней вставлено T(x\_1,x\_2).\par
Кандидаты: x\_1,x\_2. Чтение не установлено в дипломатическом исходнике.\par
Исходный PNG 2408×3480; область x=1195–1510, y=3185–3310 пикселей. Просмотрено при 6× и 10×. Верхняя вставка T(x\_1,x\_2) читается; первоначальная перечёркнутая запись приведена лишь как кандидат.\par
Повторная сверка: Повторно просмотрены исходник и увеличения 3×, 5× и 10×. Исходные наблюдения, варианты чтения, зачёркивания и неуверенности сохранены. Математические проверки отделены от буквального чтения; частичные рисунки не утверждают неустановленных соединений или направлений.\par\medskip
\subsection*{\texttt{GELFAND-P001473}}
\par\noindent\leavevmode\phantomsection\label{GELFAND-0159-U601}\label{GELFAND-P001473:mid-cancelled-after-c}\textbf{\texttt{GELFAND-0159-U601}} (mid-cancelled-after-c).\nopagebreak[4]\par
Справа от C в строке о пространстве связностей находится небольшая густо зачёркнутая запись; под штрихами различимы лишь фрагменты букв.\par
Кандидаты: не записаны. Чтение не установлено в дипломатическом исходнике.\par
Источник 2408×3480 px; область x=1390–1570, y=1230–1350, отсчёт от левого верхнего угла. Осмотрено на оригинале, в контекстном увеличении 3× и отдельном увеличении 8×; начертание до зачёркивания надёжно не читается. Вставка «неприводимых» и нижняя приписка «(эрмитовых)» прочитаны и сохранены.\par
Повторная сверка: Повторно просмотрены исходник и увеличения 3×, 5× и 10×. Исходные наблюдения, варианты чтения, зачёркивания и неуверенности сохранены. Математические проверки отделены от буквального чтения; частичные рисунки не утверждают неустановленных соединений или направлений.\par\medskip
\subsection*{\texttt{GELFAND-P001474}}
\par\noindent\leavevmode\phantomsection\label{GELFAND-0159-U602}\label{GELFAND-P001474:upper-class-symbol}\textbf{\texttt{GELFAND-0159-U602}} (upper-class-symbol).\nopagebreak[4]\par
В пункте 3, после слов «класс для», буква перед нижним индексом 2 густо обведена; видны дуга и внутреннее пересечение. В следующей строке после «для» читается c₂.\par
Кандидаты: c\_2; e\_2. Чтение не установлено в дипломатическом исходнике.\par
Источник 2408×3480 px; область x=1420–1545, y=515–633. Осмотрено на оригинале, в контексте и отдельном увеличении 8×. Повтор в следующей строке не использован для молчаливого исправления первого знака.\par
Повторная сверка: Повторно просмотрены исходник и увеличения 3×, 5× и 10×. Исходные наблюдения, варианты чтения, зачёркивания и неуверенности сохранены. Математические проверки отделены от буквального чтения; частичные рисунки не утверждают неустановленных соединений или направлений.\par\medskip
\par\noindent\leavevmode\phantomsection\label{GELFAND-0159-U603}\label{GELFAND-P001474:mid-cancelled-notation}\textbf{\texttt{GELFAND-0159-U603}} (mid-cancelled-notation).\nopagebreak[4]\par
Перед незачёркнутым знаком связности с тильдой в пункте (9) перечёркнуто обозначение: большая A-подобная буква с нижним знаком, напоминающим набла. Поверх проведены несколько диагональных штрихов.\par
Кандидаты: A\_\textbackslash{}nabla; A\_\{\textbackslash{}widetilde\textbackslash{}nabla\}. Чтение не установлено в дипломатическом исходнике.\par
Источник 2408×3480 px; область x=400–615, y=1210–1410. Осмотрено на оригинале, в контексте и увеличении 7×. Зачёркнутое обозначение сохранено как неуверенное чтение; наличие тильды в его индексе установить нельзя.\par
Повторная сверка: Повторно просмотрены исходник и увеличения 3×, 5× и 10×. Исходные наблюдения, варианты чтения, зачёркивания и неуверенности сохранены. Математические проверки отделены от буквального чтения; частичные рисунки не утверждают неустановленных соединений или направлений.\par\medskip
\par\noindent\leavevmode\phantomsection\label{GELFAND-0159-U604}\label{GELFAND-P001474:mid-remark-initials}\textbf{\texttt{GELFAND-0159-U604}} (mid-remark-initials).\nopagebreak[4]\par
В мелкой приписке после «Замеч.» написаны инициалы И., затем густой знак с двумя высокими штрихами и поперечным соединением, затем Г.\par
Кандидаты: И. М. Г.; И. Н. Г.. Чтение не установлено в дипломатическом исходнике.\par
Источник 2408×3480 px; область x=393–610, y=1630–1698. Осмотрено на оригинале, в контекстных увеличениях 4× и отдельно 8×; средняя буква остаётся неоднозначной.\par
Повторная сверка: Повторно просмотрены исходник и увеличения 3×, 5× и 10×. Исходные наблюдения, варианты чтения, зачёркивания и неуверенности сохранены. Математические проверки отделены от буквального чтения; частичные рисунки не утверждают неустановленных соединений или направлений.\par\medskip
\par\noindent\leavevmode\phantomsection\label{GELFAND-0159-U605}\label{GELFAND-P001474:mid-cancelled-line}\textbf{\texttt{GELFAND-0159-U605}} (mid-cancelled-line).\nopagebreak[4]\par
Между «Посмотрим на критич. точки» и «Это связности с кривизной 0» короткая запись зачёркнута несколькими толстыми почти горизонтальными линиями. У первого знака остаются верхняя дуга и нижняя петля; фрагменты остальных букв видны над и под штрихами, но не складываются в надёжное чтение.\par
Кандидаты: не записаны. Чтение не установлено в дипломатическом исходнике.\par
Источник 2408×3480 px; вся запись x=340–1000, y=1760–1870. Осмотрено на оригинале, в контексте, увеличении 5× и двумя перекрывающимися увеличениями 10×. Зачёркнутая запись локализована и отмечена в транскрипции без восстановления слов по смыслу.\par
Повторная сверка: Повторно просмотрены исходник и увеличения 3×, 5× и 10×. Исходные наблюдения, варианты чтения, зачёркивания и неуверенности сохранены. Математические проверки отделены от буквального чтения; частичные рисунки не утверждают неустановленных соединений или направлений.\par\medskip
\par\noindent\leavevmode\phantomsection\label{GELFAND-0159-U606}\label{GELFAND-P001474:lower-equation-label}\textbf{\texttt{GELFAND-0159-U606}} (lower-equation-label).\nopagebreak[4]\par
Слева от формулы для ∇ad номер в скобках имеет читаемую первую 1 и густо переправленную вторую цифру.\par
Кандидаты: 11. Чтение не установлено в дипломатическом исходнике.\par
Источник 2408×3480 px; область x=145–325, y=2455–2610. Осмотрено на оригинале, в контексте и увеличении 8×. Кандидат 11 сохранён как неуверенное чтение; исходная переправленная цифра не восстановлена.\par
Повторная сверка: Повторно просмотрены исходник и увеличения 3×, 5× и 10×. Исходные наблюдения, варианты чтения, зачёркивания и неуверенности сохранены. Математические проверки отделены от буквального чтения; частичные рисунки не утверждают неустановленных соединений или направлений.\par\medskip
\par\noindent\leavevmode\phantomsection\label{GELFAND-0159-U607}\label{GELFAND-P001474:lower-m-x-notation}\textbf{\texttt{GELFAND-0159-U607}} (lower-m-x-notation).\nopagebreak[4]\par
В формуле (10) непосредственно читается Ω¹(M). В следующей формуле, с переправленным номером — кандидатом (11), непосредственно читаются Ω⁰(X) и Ω¹(X). В числителе и знаменателе формулы (13) также написаны соответственно Ω¹(X) и Ω⁰(X).\par
Кандидаты: Ω¹(M) в (10); Ω⁰(X), Ω¹(X) в (11); Ω¹(X), Ω⁰(X) в (13).. Чтение не установлено в дипломатическом исходнике.\par
Сверка обозначений по источнику 2408×3480 px: формула (10), x=780–1870, y=2260–2445; формула (11), x=440–1850, y=2460–2650; формула (13), x=365–1370, y=2880–3160. Все три области непосредственно повторно осмотрены в увеличении 3×. Буквы M и X прочитаны уверенно. Это запись внешнего контроля обозначений: X сохранено дословно; новое равенство или содержательная поправка к формулам не вводится.\par
Повторная сверка: Повторно просмотрены исходник и увеличения 3×, 5× и 10×. Исходные наблюдения, варианты чтения, зачёркивания и неуверенности сохранены. Математические проверки отделены от буквального чтения; частичные рисунки не утверждают неустановленных соединений или направлений.\par\medskip
